\begin{abstract}
A common faliure mode in assistant-like AI usecases is ambiguous requests from the user. Assistants usually do not understand the request and misinterpret it or even answer incorrectly. We study \emph{intent stabilization} by sampling structured intent from user requests. We then repair the model understanding path when it is unstable and clarify when there is still uncertainty left. The method is inference time and is not intrusive, so it works with any model. On InfoQuest and ClarifyBench test splits with Qwen3-30B and a Llama-3.3-70b judge model, we find that all-stage resampling matches generic clarification on InfoQuest task-success rate (0.240) and slightly improves upon targeted clarification (0.235), but falls behind the highest scoring ClarifyBench clarification baseline (0.500 vs. 0.544). The K-sensitivity runs show that more samples are not monotonically better: \(K=7\) gives the best dev task-success rate (0.240) but costs 26.95 task-model calls per example, while \(K=1\) is close (0.230) at 6.07 calls. We find that structured resampling is useful for diagnosing unstable intent, but simple clarification is still a strong quality-cost baseline.
\end{abstract}